In February 2024, the Securities and Exchange Commission shortened the initial
Schedule~13D filing deadline from ten calendar days to five business days and
declined to extend beneficial-ownership reporting to certain cash-settled
derivatives. This paper examines whether the reform changed activist campaign
formation, the composition of activist positions, and the timing of takeover
gains. The empirical design uses a predetermined target-level exposure measure:
the number of trading days required to acquire an additional five percentage
points of ownership at a participation cap. This measure captures how tightly
the shorter deadline
constrains undisclosed post-threshold accumulation. I use it to test whether
campaign entry, target composition, stakes at filing, and Item~6 derivative
disclosure respond more strongly for more exposed targets. I also decompose
target gains into the filing return, the post-filing run-up, and the bid markup,
holding the trigger-to-bid horizon fixed. A parsimonious model of accumulation
over a filing window derives the exposure measure and identifies the conditions
under which a shorter deadline reduces the stake acquired before disclosure.
The design reports pre-specified minimum detectable effects and uses the SEC's
published tables as a pre-reform benchmark. It uses realised behaviour to
evaluate the SEC's projected \$810 million annual loss in shareholder value.
